\section{Future Work}
The work presented in this thesis advanced a compact wireless strain-sensing system from concept to a verified prototype, but several avenues remain open both to further advance the implementation and to make it deployable as a finished product. 
This section outlines the main directions for further development, covering the sensor, the hardware, the wireless communication, the firmware, and the long-term qualification of the design.
 
\subsection{Full Bridge and Reference Arm Placement}
The present sensor system uses the Wheatstone half-bridge of \Cref{Wheatstone Half-Bridge Design}, in which two active arms are mounted on the strained surface and two passive reference arms are held unstrained for temperature compensation. Possible solutions include mounting the reference arms on a neighboring, unstrained region of the metal chassis, suspending them in free air, or making use of the nut's surface area to realize a full bridge configuration instead. The first two solutions are more complex to mount, as they require routing the reference arms away from the nut face rather than keeping the whole sensor assembly self-contained, which makes the third option, realizing a full bridge directly on the nut surface, the more attractive path.

As shown in \Cref{sec:Wheatstone_theory}, replacing the two passive arms with active elements responding in the opposite direction forms a full bridge, which ideally doubles the output swing once more and linearizes the output in the relative resistance change. Moving to a full bridge would therefore roughly double the sensitivity and improve the rejection of common-mode and temperature effects, at the cost of two additional active elements and the more demanding matching this requires.

A practical way to realize this is to mount all four arms on the strained surface but orient the two reference arms at \SI{90}{\degree} to the strain axis. Aligned transversely, these arms respond primarily to the Poisson strain, which is of opposite sign and smaller magnitude than the axial strain seen by the active arms, so they contribute a partial active response in the opposing direction while remaining far less sensitive than the axial arms~\cite{hbm_wheatstone_bridge}. The advantage of this arrangement is mainly thermal, since with all four arms printed together and bonded to the same surface they are held at the same temperature and experience the same thermal environment, which directly addresses the matching and equal-temperature assumptions that \Cref{Wheatstone Half-Bridge Design} identifies as the limiting factors in the compensation.

\subsection{Print-Process Control and Sensor Reproducibility}\label{sec:process_control}
The trade-off examined in \Cref{sec:tradeoff} traced the print-to-print variation of the sensor to small differences in the print, most directly the nozzle-to-substrate spacing, whose effect is amplified at the near-separation operating point at which the sensor is most sensitive. The sensors in this work were produced on the printer available for the project, whose control over the extrusion rate and this spacing sets a floor on how closely two elements can be made to match. A more capable machine, or one purpose-built for this deposition, with closed-loop extrusion control and a temperature-regulated build environment, would be expected to narrow this variation at its source, improving both the reproducibility of the print optimization and the matching of the bridge arms, and with it the common-mode rejection on which the temperature compensation depends. Such an improvement would also sharpen the smaller optimization effects that the present print-to-print spread tends to obscure. Beyond the manufacturing benefit set out in \Cref{sec:tradeoff}, where reproducible printing would allow a single type-calibration to replace per-unit calibration, tighter process control is therefore the most direct route to addressing the sensitivity-reproducibility trade-off at its origin.
 
\subsection{Sensor Geometry and Mechanical Repeatability}
Two further refinements concern the printed element itself. The serpentine currently spans \SI{25}{\milli\meter} of the \SI{30}{\milli\meter} available on the nut face, so extending it to the full width offers a direct route to additional sensitivity without altering the sensing principle. Separately, the hysteresis under repeated loading and the initial settling observed in the rest-stability measurements were attributed to viscoelastic relaxation of the substrate and to the gradual settling of the asperity contacts, so a stiffer or less viscoelastic substrate, or a mechanical break-in procedure applied before calibration, would be worth investigating as a means of improving the short-term repeatability. How these properties evolve over years of service is left to the accelerated life testing described in \Cref{sec:improved_test}.

A complementary refinement would target the substrate rather than the machine. In the present process the conductive element is deposited onto a printed substrate whose own thickness varies from one print to the next, and this variation carries directly into the effective nozzle-to-substrate spacing, adding to the same z-height uncertainty to which the near-separation operating point makes the sensor so sensitive. Printing the element instead onto a pre-manufactured TPU sheet of precise and constant thickness would remove this contribution, since a substrate of known and repeatable thickness fixes the height of the printing surface between prints and leaves the machine's own positioning as the only remaining source of z-height variation. This would be expected to improve the print-to-print consistency of the resting contact, and with it the reproducibility of the sensor, at the operating point where such variation has the greatest effect.

\subsection{Field Testing and Wireless Communication}
The module was validated in the laboratory and climate chamber, but not in its intended environment. A natural next step is to deploy it on an operational offshore wind turbine, where the metallic nacelle and tower, the rotating machinery, and the long internal signal paths create propagation conditions that cannot be reproduced on the bench. Such a deployment would establish the true range and link margin and reveal environmental effects, such as vibration and electromagnetic interference, that laboratory testing could not capture.
 
Should this testing show that the \SI{2.4}{\giga\hertz} link does not provide sufficient range or margin through the turbine structure, a study of sub-\si{\giga\hertz} implementations would be worthwhile. A sub-\si{\giga\hertz} signal is less attenuated by the metallic structure and less affected by interference from WiFi and Bluetooth, yielding additional link margin for the same transmit power~\cite{silabs_subghz_priorities}. This margin can be spent in either of two ways, extending the range or, alternatively, lowering the transmit power while still closing the link. Since the radio transmission is one of the largest contributors to the energy consumed in each measurement cycle, a lower transmit power reduces the energy per transmission and thereby extends the battery life. Such a change does not require abandoning Zigbee, as it is one of the few protocols available in both \SI{2.4}{\giga\hertz} and sub-\si{\giga\hertz} versions. Alternatively, a dedicated sub-\si{\giga\hertz} protocol such as Z-Wave could be considered, at the cost of a smaller and more expensive device ecosystem~\cite{homey_zigbee_vs_zwave}.

\subsection{Sensor Commissioning}
In the current implementation, each module is identified manually by assigning it a name in Home Assistant as it joins the coordinator. This step is workable for a few modules, but it does not scale well to a larger deployment and is prone to human error, as the installer must correctly match each joining device to its intended name and physical location. A useful extension would be to streamline this commissioning step, for example by labeling each module with a bar- or QR code encoding its identifier, so that it can be scanned and registered to a known position during installation rather than named manually. This would make the mapping between modules and locations faster, less error-prone, and better suited to deployments involving many modules.

\subsection{Temperature Compensation}
The characterization of the on-die temperature sensor presented in \Cref{sec:temp_sensor} established that it provides sufficient precision to serve as a reference for compensating the residual temperature sensitivity introduced by imbalance in the Wheatstone bridge. A natural extension is therefore the derivation and integration of a compensation model that expresses the measured bridge output as a function of die temperature, allowing the residual thermal drift to be removed in firmware without modification to the existing hardware.

\subsection{Battery Monitoring}
The current design does not measure its own battery voltage, so the state of the battery is not available in the user interface. For a system intended to operate unattended for years, this is a significant limitation, as there is no way to detect a weakening battery before it fails. The measurements already show a strong temperature dependence of the battery voltage (\Cref{fig:FS_bat}), which makes such monitoring especially valuable, since a module may approach its voltage limit at low temperature while still appearing healthy under normal conditions. A useful extension is therefore to add battery-voltage sensing to the module, allowing the battery state of each module to be reported alongside its measurements and surfaced in the interface.

\subsection{Encapsulation Validation}
The present encapsulation encloses the PCB and sensor unit in a silicone compound (\Cref{sec:encapsulation}). This was found to attenuate the RF link only slightly (\Cref{sec:signal_strength}) and performed well overall, but its protective properties were not formally validated. Future work should therefore characterize the encapsulation more thoroughly, including extended environmental testing, and verification of resistance to oil, grease, and water ingress. In parallel, alternative encapsulants could be investigated, with the aim of finding a solution that is both more economical and less attenuating to the RF signal than the current silicone compound.

\subsection{Improved Test Methodology and Long-term Qualification} \label{sec:improved_test}
The characterization in this work was performed with the available laboratory equipment, which limited the precision of the applied mechanical load and the simultaneous control of strain and temperature. A more rigorous characterization would use a calibrated high-force hydraulic load frame, such as one of the systems in Instron's Industrial Series~\cite{instron_industrial}, whose force capacity spans the range required to apply the bolt's service preload, in which a dedicated fixture places the bolt under a controlled and repeatable tensile load that is then released, while strain and temperature are measured simultaneously. 
This would allow the sensor response, during both loading and relaxation, to be related directly to a known applied load across the operating range.
Straining the bolt in this way also exercises the sensor in its deployed, surface-bonded configuration rather than clamped between two grips, which reproduces the compressive loading the deployed sensor experiences and removes the free-span buckling noted in \Cref{sec:reliability}.
Beyond the sensor itself, the durability of the system has not been established. Accelerated life testing and ingress testing of the encapsulation would quantify how the protection and the sensor response evolve over time and under repeated environmental cycling, while a field trial in the intended environment would finally confirm that the system performs as expected under real operating conditions over an extended period. Together, these steps would move the system from a verified prototype towards a validated, qualified, deployable product.